\clearpage
\fancyhead[L]{\small GAUSS TO VON LINDENAU.}
\SourceLeaf{310}\begin{center}
GAUSS TO VON LINDENAU.
\end{center}

\hrule
\begin{center}
\fontsize{10}{12}\selectfont Monatliche Correspondenz. Vol. XVIII. Pp. 173..175. August 1808.
\end{center}
\hrule

\begin{flushright}
Göttingen, 6 August 1808.
\end{flushright}

\par\indent Since my last letter I have continued observing the new planets whenever circumstances allowed, and I am glad to share what I have found so far. Until now, \emph{Juno}, \emph{Vesta}, and \emph{Pallas} could only be observed with the ring micrometer. I have made no further observations of Vesta since 13 July, however, as it can now be observed on the meridian. Professor \textsc{Harding} also succeeded in observing \emph{Ceres} with the quadrant. I will send you that observation later [handwritten remark: July 25. $13^{\text{u}}\,14^{\text{m}}\,29^{\text{s}}$ \quad $322^\circ\,8'\,57''6$ \quad $27^\circ\,52'\,11''5$]

\par\indent Unfortunately, \emph{Pallas} now stands too high for our observing site all night. I was able to obtain only a few observations, and their declinations are uncertain.

\par\indent Mr. \textsc{Harding} made the observation of \emph{Vesta} on 2 July and the one of \emph{Pallas} on 15 July. I made all the others. For \emph{Pallas}, the ephemeris has an error of $-4'$ in right ascension and gives the declination to within $1'$. For \emph{Juno}, the difference in right ascension has now grown to $13'$. I consider all four of the latest observations of \emph{Juno} very good, insofar as the comparison stars were well observed in the Hist. Cél.

\Needspace{13\baselineskip}
\begin{center}
\emph{Observations of Pallas.}
\end{center}


\begin{center}\fontsize{10}{12}\selectfont\setlength{\tabcolsep}{3pt}\renewcommand{\arraystretch}{1.15}\begin{tabular}{lrrr}
1808 & \shortstack{Mean time\\at Göttingen} & \shortstack{Apparent\\R.A.} & \shortstack{Apparent\\declination}\\\hline
July 13 & $12^{\text{u}}\;3^{\text{m}}\;17^{\text{s}}$ & $300^\circ\;21'\;49''$ & $18^\circ\;41'\;44''$ N.\\
15 & $12\;1\;28$ & $299\;57\;3$ & $18\;32\;40$\\
25 & $10\;9\;44$ & $297\;57\;37$ & $17\;40\;10$\\
\end{tabular}\end{center}


\Needspace{13\baselineskip}
\begin{center}
\emph{Observations of Juno.}
\end{center}


\begin{center}\fontsize{10}{12}\selectfont\setlength{\tabcolsep}{3pt}\renewcommand{\arraystretch}{1.15}\begin{tabular}{lrrr}
1808 & \shortstack{Mean time\\at Göttingen} & \shortstack{Apparent\\R.A.} & \shortstack{Apparent\\declination}\\\hline
June 20 & $11^{\text{u}}\;49^{\text{m}}\;0^{\text{s}}$ & $315^\circ\;29'\;34''$ & $2^\circ\;16'\;23''$ S.\\
22 & $12\;0\;45$ & $315\;23\;1$ & $2\;14\;29$\\
July 6 & $12\;42\;23$ & $314\;1\;0$ & $2\;19\;7$\\
30 & $10\;51\;17$ & $309\;38\;9$ & $3\;57\;13$\\
31 & $11\;40\;57$ & $309\;24\;53$ & $4\;3\;48$\\
Aug. 4 & $10\;42\;53$ & $308\;33\;44$ & $4\;30\;32$\\
5 & $10\;53\;46$ & $308\;20\;39$ & $4\;37\;36$\\
\end{tabular}\end{center}


\par\indent [The original also includes observations of \emph{Vesta} from June 22 and July 1, 2, and 13. They also appear in the notice from the Gött. gel. Anz. of 1808 Nov. 10, reprinted below on p. 309.]

\par\indent I have already calculated the opposition of \emph{Juno} levi calamo from my observations and made a preliminary adjustment to the elements. A very small change is enough to represent all four oppositions observed so far very accurately with a pure ellipse. It consists mainly in a slight decrease in the longitude of perihelion and a slight increase in the mean motion. More on this soon.

\par\indent The King has given our observatory several astronomical instruments. A Gregorian telescope by \textsc{Chevalier} and a marine chronometer by \textsc{Louis Berthoud} have already ar\SourceLeaf{311}rived. As far as we can tell from our experience so far, the chronometer runs fairly steadily. Of course, its rate has so far been observed only while the instrument was kept stationary. We will soon try to determine a longitude with it. [This notice about the instruments is also printed in the Gött. gel. Anzeigen, issue 127, page 1265, 1808 August 8, with the addition “We may also expect the other instruments, which have been in Paris until now, to arrive soon.” \textsc{Schering}.]

